\subsection{Ejercicio 5.1: implementación de los operadores y generación de matrices}
\label{sec:ejercicio51}

En este ejercicio se implementaron los operadores $\varphi$, $\theta$ y
$\gamma$ descritos por Alba-Cabrera et al.~\cite{alba2016}. Su propósito
es construir matrices básicas con dimensiones controladas y con un número
de testores típicos que puede calcularse de antemano. Esto permite preparar
experimentos cuya respuesta esperada es conocida, en lugar de elegir
matrices aleatorias sin una referencia para comprobar los resultados.

Se conserva el procedimiento de las Tablas 5 y 6 del artículo, pero se
utilizan las matrices $A$ y $B$ de nuestros ejercicios 2.1 y 4.1.
Por tanto, se reproducen las familias de operaciones y los valores de $N$,
no los valores numéricos de las matrices originales del artículo.
La generación de matrices corresponde al punto 5.1; sus tiempos de
resolución con YYC y BT corresponden al punto 5.2.

\subsubsection*{Definición de los operadores}

Sean $U$ una matriz básica de $m\times n$ y $V$ una matriz básica de
$p\times q$. Al unirlas, las columnas de $V$ ocupan posiciones nuevas,
desde $n+1$ hasta $n+q$. No se identifican con las columnas de $U$,
aunque en sus archivos originales ambas matrices comiencen por $x_1$.

\paragraph{Operador $\theta$: todas las parejas de filas.}
Cada fila de $U$ se concatena horizontalmente con cada fila de $V$:
\[
\theta(U,V)=
\begin{pmatrix}
u_1\mid v_1\\
u_1\mid v_2\\
\vdots\\
u_1\mid v_p\\
u_2\mid v_1\\
\vdots\\
u_m\mid v_p
\end{pmatrix}.
\]
La notación $u_i\mid v_j$ significa escribir primero los $n$ valores
de $u_i$ y después los $q$ valores de $v_j$; la barra sólo separa
visualmente los bloques. El resultado tiene $mp$ filas y $n+q$
columnas. No requiere que las entradas tengan el mismo número de filas.

El procedimiento computacional recorre las filas de $U$ con un ciclo
externo y las de $V$ con un ciclo interno. Así produce, en ese orden,
$(u_1\mid v_1),\ldots,(u_1\mid v_p),(u_2\mid v_1),\ldots$.
No se utiliza la multiplicación matricial usual ni el producto de
Kronecker: la operación definida aquí es una concatenación para cada
pareja de filas.

\paragraph{Operador $\varphi$: concatenación fila por fila.}
Requiere $m=p$. Une la primera fila de $U$ con la primera de $V$,
la segunda con la segunda, y así sucesivamente:
\[
\varphi(U,V)=[\,U\mid V\,]
=\begin{pmatrix}
u_1\mid v_1\\
u_2\mid v_2\\
\vdots\\
u_m\mid v_m
\end{pmatrix}.
\]
El resultado tiene $m$ filas y $n+q$ columnas. A diferencia de
$\theta$, no genera todas las parejas posibles. Para obtener
$\varphi^N(U)$ se colocan $N$ copias de $U$ horizontalmente:
\[
\varphi^N(U)=[\,\underbrace{U\mid U\mid\cdots\mid U}_{N\text{ copias}}\,].
\]
Así se conservan las $m$ filas y se obtienen $Nn$ columnas.

\paragraph{Operador $\gamma$: bloques diagonales.}
Coloca $U$ arriba a la izquierda y $V$ abajo a la derecha, y completa
con ceros los otros dos bloques:
\[
\gamma(U,V)=
\begin{pmatrix}
U & 0_{m\times q}\\
0_{p\times n} & V
\end{pmatrix}.
\]
Su tamaño es $(m+p)\times(n+q)$. Para $\gamma^N(U)$ se colocan $N$
copias de $U$ en la diagonal y ceros fuera de esos bloques. En este
caso se obtienen $Nm$ filas y $Nn$ columnas.

En la notación del artículo, $N$ cuenta las copias de la matriz,
no el número de llamadas adicionales a la función. Por ello,
$\varphi^1(U)=\gamma^1(U)=U$ y
$\gamma^2(U)=\gamma(U,U)$.

\subsubsection*{Ejemplo para distinguir las tres construcciones}

Considérese la matriz pequeña $P=I_2$, independiente de las matrices
$A$ y $B$ del proyecto:
\[
P=\begin{pmatrix}1&0\\0&1\end{pmatrix}.
\]
Aplicando cada operador a dos copias de $P$, se obtiene:
\[
\theta(P,P)=\begin{pmatrix}
1&0&1&0\\1&0&0&1\\0&1&1&0\\0&1&0&1
\end{pmatrix},\qquad
\varphi(P,P)=\begin{pmatrix}
1&0&1&0\\0&1&0&1
\end{pmatrix},\qquad
\gamma(P,P)=\begin{pmatrix}
1&0&0&0\\0&1&0&0\\0&0&1&0\\0&0&0&1
\end{pmatrix}.
\]
Las tres matrices tienen cuatro columnas, pero se construyen de
manera distinta: $\theta$ combina todas las parejas, $\varphi$ une
filas correspondientes y $\gamma$ separa las copias en bloques
diagonales. En este ejemplo tienen, respectivamente, 4, 2 y 4 filas.

\subsubsection*{Conservación de la basicidad}

Los operadores conservan la condición de matriz básica cuando sus
entradas son básicas y no contienen filas nulas~\cite{alba2016}. Esta propiedad puede entenderse
comparando los bloques de las filas:
\begin{itemize}[nosep]
\item En $\theta(U,V)$, si una fila fuera subfila de otra, su bloque
de $U$ tendría que estar contenido en el otro bloque de $U$, y lo
mismo ocurriría con los bloques de $V$. La basicidad obliga a que
ambos bloques sean iguales, por lo que no existe una inclusión estricta.
Además, parejas distintas no producen filas repetidas.
\item En $\varphi(U,V)$, dos filas diferentes ya son incomparables
en el bloque correspondiente a $U$. Agregar el bloque de $V$ no
elimina esa diferencia ni puede convertirlas en filas repetidas.
\item En $\gamma(U,V)$, dentro de cada bloque se conserva la
incomparabilidad original. Entre bloques distintos, cada fila
tiene algún 1 en posiciones donde la otra tiene sólo ceros.
\end{itemize}
Las filas tampoco son nulas, porque cada una conserva una fila
no nula de alguna matriz de entrada. Por tanto, no es necesario
reducir ni reordenar los resultados antes de usarlos en los experimentos.
Esta afirmación se comprobó también con la rutina de matriz básica.

\subsubsection*{Implementación computacional}

La implementación está escrita en C++17 y utiliza la representación
\texttt{BooleanMatrix}. Cada fila es un \texttt{ColumnSet}, que
almacena las posiciones de sus unos. Las funciones públicas se
declaran en \texttt{include/testores/domain/matrix\_operators.hpp}
y se implementan en \texttt{src/domain/matrix\_operators.cpp}.

\begin{center}
\begin{tabular}{l p{9cm}}
\toprule
Función & Responsabilidad\\
\midrule
\texttt{theta(U,V)} & Generar todas las concatenaciones $u_i\mid v_j$.\\
\texttt{phi(U,V)} & Concatenar las filas correspondientes, comprobando que haya igual cantidad de filas.\\
\texttt{phiPower(U,N)} & Construir directamente las $N$ copias horizontales de $U$.\\
\texttt{gamma(U,V)} & Construir los dos bloques diagonales y sus ceros.\\
\texttt{gammaPower(U,N)} & Construir directamente $N$ bloques diagonales de $U$.\\
\bottomrule
\end{tabular}
\end{center}

Para concatenar un bloque a la derecha, las posiciones de sus unos
se desplazan por el ancho del bloque izquierdo. Por ejemplo, si
éste tiene siete columnas, la primera columna del bloque derecho
ocupa la posición $x_8$. En las potencias de $\gamma$, las posiciones
fuera del bloque activo permanecen en cero.

Las funciones rechazan entradas vacías, cero copias y resultados
que excedan los límites de 100\,000 filas o 512 columnas. La
validación de dimensiones se realiza antes de reservar la memoria.
La interfaz de comandos exige entradas básicas y no las reduce
automáticamente, para conservar el orden de generación.

Las pruebas automáticas incluyen ejemplos con resultados conocidos,
dimensiones incompatibles, límites de recursos y desplazamientos
que atraviesan posiciones de 64 bits. En ejemplos pequeños se
comprueba además que YYC y BT recuperen los testores típicos esperados.

\subsubsection*{Construcción de la matriz base del experimento}

Los datos de partida de este proyecto son:
\begin{center}
\begin{tabular}{c r r r l}
\toprule
Matriz & Filas & Columnas & $|\Psi^*|$ & Tamaños de los testores típicos\\
\midrule
$A$ & 6 & 7 & 8 & 5 de tamaño 2 y 3 de tamaño 3\\
$B$ & 8 & 7 & 6 & 4 de tamaño 3 y 2 de tamaño 4\\
\bottomrule
\end{tabular}
\end{center}

Como $A$ y $B$ tienen distinta cantidad de filas, no se aplica
$\varphi(A,B)$ directamente. Siguiendo las Tablas 5 y 6 del artículo,
primero se construye
\[
C=\theta(A,B),\qquad
\text{filas}(C)=6\cdot8=48,\qquad
\text{columnas}(C)=7+7=14.
\]
Las primeras ocho filas son $a_1\mid b_1,\ldots,a_1\mid b_8$;
las siguientes ocho corresponden a $a_2$, y así sucesivamente.
Por ejemplo, la primera fila se obtiene de
\[
a_1=(1,0,1,0,1,0,1),\qquad b_1=(1,0,1,0,0,0,0),
\]
y queda
\[
c_1=(1,0,1,0,1,0,1\mid1,0,1,0,0,0,0).
\]
El bloque de $A$ ocupa $x_1,\ldots,x_7$ y el de $B$ ocupa
$x_8,\ldots,x_{14}$. Esta separación es importante al contar
los testores típicos.

\subsubsection*{Cálculo previo del número de testores típicos}

Las proposiciones del artículo permiten conocer la respuesta
esperada antes de resolver las matrices grandes. En $\theta(U,V)$,
los testores típicos son los de $U$ o los de $V$, cada grupo dentro
de su propio bloque de columnas. Si un conjunto no cubriera ni
$U$ ni $V$, existiría una pareja de filas cuya concatenación
tampoco quedaría cubierta. Por tanto:
\[
|\Psi^*(\theta(U,V))|=|\Psi^*(U)|+|\Psi^*(V)|,
\qquad |\Psi^*(C)|=8+6=14.
\]
De los 14 testores de $C$, cinco tienen tamaño 2, siete tienen
tamaño 3 y dos tienen tamaño 4. Los siete de tamaño 3 provienen
de los tres de $A$ y los cuatro de $B$.

En $\varphi^N(C)$ cada columna original aparece en $N$ copias
idénticas. Un testor típico de tamaño $k$ origina $N^k$
selecciones: se elige una copia para cada una de sus $k$
características. No se necesitan dos copias de la misma
característica, porque una sería redundante. En consecuencia:
\[
|\Psi^*(\varphi^N(C))|
=\sum_{T\in\Psi^*(C)}N^{|T|}
=5N^2+7N^3+2N^4.
\]
Por ejemplo, para $N=2$ se esperan
$5\cdot4+7\cdot8+2\cdot16=108$ testores típicos.

En $\gamma(U,V)$ se deben cubrir ambos bloques, de modo que
cada testor típico resulta de combinar uno de $U$ con uno de
$V$, desplazando las columnas del segundo bloque:
\[
|\Psi^*(\gamma(U,V))|=|\Psi^*(U)|\,|\Psi^*(V)|.
\]
Al repetir $C$ en $N$ bloques diagonales, hay 14 elecciones
por bloque y las combinaciones son independientes:
\[
|\Psi^*(\gamma^N(C))|=14^N.
\]
Por ejemplo, $\gamma^2(C)$ debe tener $14^2=196$ testores
típicos. Estos conteos son resultados teóricos utilizados para
validar las ejecuciones, no tiempos de los algoritmos.

\subsubsection*{Matrices generadas para las Tablas 5 y 6}

Para reproducir la familia de la Tabla 5 se generaron
$\varphi^N(C)$ con $N=1,2,3,4,5$. Para la Tabla 6 se generaron
$\gamma^N(C)$ con $N=1,2,3,4$. El caso $N=1$ es $C$ en ambas
familias, por lo que sólo se necesitan ocho archivos distintos.

\begin{table}[htbp]
\centering
\small
% Generado por generar_ejercicio_5_1.ps1: conteos teóricos, no tiempos.
\begin{tabular}{c l r r r}
\toprule
$N$ & Matriz & Filas & Columnas & $|\Psi^*|$ teórico \\
\midrule
1 & $\varphi^{1}(C)$ & 48 & 14 & 14 \\
2 & $\varphi^{2}(C)$ & 48 & 28 & 108 \\
3 & $\varphi^{3}(C)$ & 48 & 42 & 396 \\
4 & $\varphi^{4}(C)$ & 48 & 56 & 1040 \\
5 & $\varphi^{5}(C)$ & 48 & 70 & 2250 \\
\bottomrule
\end{tabular}

\caption{Familia de la Tabla 5 del artículo, adaptada a nuestras
matrices $A$ y $B$. Se mantienen 48 filas y aumentan las columnas.}
\label{tab:generacionphi51}
\end{table}

\begin{table}[htbp]
\centering
\small
% Generado por generar_ejercicio_5_1.ps1: conteos teóricos, no tiempos.
\begin{tabular}{c l r r r}
\toprule
$N$ & Matriz & Filas & Columnas & $|\Psi^*|$ teórico \\
\midrule
1 & $\gamma^{1}(C)$ & 48 & 14 & 14 \\
2 & $\gamma^{2}(C)$ & 96 & 28 & 196 \\
3 & $\gamma^{3}(C)$ & 144 & 42 & 2744 \\
4 & $\gamma^{4}(C)$ & 192 & 56 & 38416 \\
\bottomrule
\end{tabular}

\caption{Familia de la Tabla 6 del artículo, adaptada a nuestras
matrices $A$ y $B$. Crecen las dimensiones y el número de testores
típicos aumenta exponencialmente.}
\label{tab:generaciongamma51}
\end{table}

En la primera familia, las dimensiones son $48\times14N$.
En la segunda son $48N\times14N$. Todas las matrices generadas
resultaron básicas: la rutina de reducción no eliminó filas
nulas, repetidas ni superfilas. Se mantuvo su orden de generación;
las copias ordenadas para comparar tiempos se tratan por separado.

\subsubsection*{Cómo reproducir la generación}

Desde la carpeta principal del proyecto, después de compilar
en Release, se puede ejecutar en PowerShell:
\begin{Verbatim}[fontsize=\small, frame=single]
.\scripts\generar_ejercicio_5_1.ps1
\end{Verbatim}
El script utiliza los archivos \texttt{data/matriz\_A.txt} y
\texttt{data/matriz\_B.txt}. Guarda una nueva serie en
\texttt{data/ej5\_1/} y las tablas y comprobaciones en
\texttt{docs/resultados/ej5\_1/}. Si esas carpetas ya existen,
se debe indicar otra salida o añadir \texttt{-Forzar} para
reemplazar esa generación. Estas matrices son idénticas a las que
guarda \texttt{testores\_experimentos} en \texttt{docs/resultados/}.

También se puede generar cada operación individualmente:
\begin{Verbatim}[fontsize=\small, frame=single]
testores theta --entrada-a data/matriz_A.txt --entrada-b data/matriz_B.txt \
               --salida ejemplo_theta.txt
testores phi     --entrada ejemplo_theta.txt --copias 2 --salida ejemplo_phi2.txt
testores gamma   --entrada ejemplo_theta.txt --copias 2 --salida ejemplo_gamma2.txt
testores ordenar --entrada ejemplo_theta.txt --salida ejemplo_theta_asc.txt
\end{Verbatim}

El script comprueba todas las filas generadas contra las
concatenaciones y bloques esperados, verifica que las matrices
sean básicas y registra sus dimensiones y hashes. Para calcular
los conteos teóricos obtiene primero con YYC los testores de
$A$ y $B$, que son matrices pequeñas. No ejecuta BT sobre las
matrices grandes ni presenta sus conteos analíticos como si
fueran resultados de una ejecución larga.

Con ello queda preparada la entrada del experimento: las
matrices se generan de forma reproducible y sus propiedades
esperadas están determinadas. La comparación de tiempos de
YYC y BT, con y sin ordenar las filas, constituye el punto 5.2.
